\section{Сравнение с предыдущей версией}

\subsection{Время и память}
Лучшее время из 5 запусков, пиковый RSS по {\ttfamily /usr/bin/time}:

\begin{longtable}{|l|r|r|r|r|r|}
\hline
\rowcolor{lightgray}
Тест & \multicolumn{2}{c|}{Старая} & \multicolumn{2}{c|}{Новая} & Ускорение\\
\rowcolor{lightgray}
 & время, с & память, МБ & время, с & память, МБ & \\
\hline
{\ttfamily mix\_100000}      & 0.071 & 6.1  & 0.042 & 4.9  & 1.7\\
{\ttfamily mix\_1000000}     & 0.893 & 31.2 & 0.537 & 19.2 & 1.7\\
{\ttfamily insert\_1000000}  & 1.357 & 80.6 & 0.769 & 47.4 & 1.8\\
{\ttfamily sorted\_1000000}  & 0.642 & 79.6 & 0.320 & 49.1 & 2.0\\
{\ttfamily saveload\_200000} & 0.867 & 35.1 & 0.429 & 23.6 & 2.0\\
\hline
\end{longtable}

Новая версия быстрее в 1.7--2 раза и тратит на 33--41\% меньше памяти на больших тестах.
Файл словаря тоже стал немного меньше (3.7 МБ против 3.8 МБ на $2\cdot10^5$ слов), так как длина
пишется одним байтом вместо двух.

\subsection{gprof}
\begin{alltt}
\$ gprof -b -p new/bpg gmon.out           # mix_1000000
  \%   cumulative   self              self     total
 time   seconds   seconds    calls  ns/call  ns/call  name
 27.79      0.10     0.10 10132604     9.87     9.87  balance(Node*)
 25.01      0.19     0.09   300389   299.73   299.73  find(Node*, char const*, int)
 22.23      0.27     0.08   499745   160.14   305.53  insert(Node*, char const*, ...)
 11.12      0.31     0.04  1000000    40.02    40.02  readWord(string const&, ...)
  5.56      0.33     0.02   199866   100.11   207.56  erase(Node*, char const*, ...)
\end{alltt}
Вызовов {\ttfamily balance} стало 10.1 млн вместо 11.8 млн. Профиль теперь покрывает 0.36~с из
0.54~с работы: основная доля времени "--- сам код словаря, а не стандартная библиотека.

На {\ttfamily saveload\_200000} один {\ttfamily Load} занимает 3.0~мс вместо 22.9~мс, то есть в 7.6 раза
быстрее, а {\ttfamily balance} при загрузке не вызывается вообще.

\subsection{callgrind}
\begin{longtable}{|l|r|r|}
\hline
\rowcolor{lightgray}
{\ttfamily mix\_100000} & Старая & Новая\\
\hline
Всего инструкций & 489.4 млн & 236.1 млн\\
{\ttfamily istringstream} + разбор числа & $\approx$ 210 млн & 0\\
{\ttfamily readWord} + {\ttfamily readNum} & "--- & 48.2 млн\\
\hline
\end{longtable}

\subsection{Куча: memcheck, dmalloc, massif}
\begin{longtable}{|l|r|r|}
\hline
\rowcolor{lightgray}
{\ttfamily mix\_100000} & Старая & Новая\\
\hline
Число выделений & 213 834 & 46 697\\
Всего выделено & 8.11 МБ & 1.93 МБ\\
Максимум одновременно (dmalloc) & 2.54 МБ & 1.41 МБ\\
Утечки & 0 & 0\\
\hline
\rowcolor{lightgray}
{\ttfamily saveload\_200000} (massif) & Старая & Новая\\
\hline
Пик кучи & 29.1 МБ & 20.1 МБ\\
Полезные данные в пике & 26.0 МБ & 15.7 МБ\\
\hline
\end{longtable}

Число выделений почти равно числу добавленных слов, как и должно быть: одна аллокация на узел.

\begin{alltt}
    MB                                                     (новая версия)
20.06^                    #
     |                    #     :    :     :    :     :     :    ::    :     @
     |                   :#   :::    @    ::   ::    ::    ::   @::   @:    :@
     |                  ::#   :::   :@   :::   ::   :::   @::  :@::  :@:  :::@
     |               ::@::#:::::@@@::@::::::@@::::::::::::@:::::@:::::@::::::@
     |          ::::@: @::#: :::@  ::@: ::::@ :::: ::::: :@::: :@:: ::@:: :::@
     | :::::::::::::@: @::#: :::@  ::@: ::::@ :::: ::::: :@::: :@:: ::@:: :::@
   0 +----------------------------------------------------------------------->Gi
     0                                                                   2.525
\end{alltt}
Ось времени у {\ttfamily massif} "--- выполненные инструкции: 2.5 млрд вместо 5.9 млрд.

\subsection{Обработка системных ошибок}
Нехватка места на диске ({\ttfamily ulimit -f 1000}, Save $10^5$ слов поверх файла в 25 байт):
\begin{longtable}{|l|l|l|}
\hline
\rowcolor{lightgray}
 & Вывод & Файл после\\
\hline
Старая & {\ttfamily ERROR: File too large} & 1 024 000 байт мусора\\
Новая & {\ttfamily ERROR: File too large} & прежние 25 байт\\
\hline
\end{longtable}

Нехватка памяти ({\ttfamily ulimit -v 30000}, тест {\ttfamily insert\_1000000}):
\begin{longtable}{|l|r|r|r|}
\hline
\rowcolor{lightgray}
 & {\ttfamily OK} & {\ttfamily Exist} & {\ttfamily ERROR: Not enough memory}\\
\hline
Старая & 265 526 & 126 198 & 608 276\\
Новая & 464 394 & 132 970 & 402 636\\
\hline
\end{longtable}
Обе версии не падают, а выводят ошибку и продолжают работу. Новая в тех же 30 МБ помещает
на 75\% больше слов.
\pagebreak
